% !PRINTED 123
\weberpage{123}

If therefore $f(b)>g$, the cut $(\mathfrak A,\mathfrak B)$ defines a number $\xi$ in the interior of the interval (1), and we can now show that
\begin{equation}
 f(\xi)=g
\tag{3}
\end{equation}
must hold.

If $f(\xi)>g$ and $f(\xi)>g'>g$, then, by the continuity of $f(x)$, we can determine an interval
\[
 \xi-h\le x\le \xi+h
\]
in which all function values $f(x)$ are greater than $g'$, and hence their \gesperrt{lower bound} is greater than $g$. But since $\xi-h$ belongs to $\mathfrak A$, the \gesperrt{lower bound} of $f(x)$ in the whole interval
\[
 a\le x\le \xi+h
\]
is also greater than $g$, while $\xi+h$ belongs to $\mathfrak B$, and this is the contradiction.

Let $f(x)$ now be an integral rational function with arbitrary complex or real coefficients, and let the variable $x$ also be allowed to take complex values.

By \S\,31, IV, if $C$ is an arbitrary positive value, the positive quantity $R$ can be so chosen that
\begin{equation}
 |f(x)|>C
\tag{4}
\end{equation}
whenever
\begin{equation}
 |x|\ge R.
\tag{5}
\end{equation}
For visualization we represent the region $(R)$ bounded by the inequality
\[
 |x|<R
\]
for the variable $x$ by a circular disc of radius $R$ in the plane of the variable $x=y+iz$.

If we take $C$ larger than any value of $|f(x)|$ in the interior of the region $(R)$, for example larger than $|f(0)|$, then $|f(x)|$ will certainly become smaller in the interior of $(R)$ than on its boundary. Since now in the whole interior of $(R)$ the function $|f(x)|$, which for abbreviation we shall now denote by $X$, is not negative, there is a \gesperrt{lower bound} $g$ for the values of $X$, and we have to prove the proposition:

4. \gesperrt{There exists a value $\xi$ of $x$ in the interior of the region $(R)$ such that}
\begin{equation}
 |f(\xi)|=g
\tag{6}
\end{equation}
%CONT>>
